\section{Chapter 4: Random Processes}

\subsection{Expected Value}

\begin{idea}
	Expectation behaves like a ``center of mass’’ for a random variable.
	Each value is weighted by how plausible it is, and the result is the balance point.
\end{idea}

For a continuous random variable with pdf $f_X(x)$:
\begin{equation}\label{eq:expectedvalue}
	\mathbb{E}[X] = \int_{-\infty}^{\infty} x\, f_X(x)\, dx.
\end{equation}

For a discrete random variable, the integral becomes a sum:
\begin{equation}\label{eq:expectedvaluediscrete}
	\mathbb{E}[X] = \sum_x x\,P(X=x).
\end{equation}

% -------------------------------------------------------------------

\subsection{Uniform Distribution}

A uniform distribution over $[a,b]$ has constant density function:
\begin{equation}\label{eq:uniformpdf}
	f_X(x) =
	\begin{cases}
		\frac{1}{b-a} & a\leq x\leq b    \\
		0             & \text{otherwise}
	\end{cases}
\end{equation}

\begin{derivation}
	Compute the expected value:
	\[
		\mathbb{E}[X]
		= \int_a^b x \cdot \frac{1}{b-a}\, dx
		= \frac{1}{b-a}\int_a^b x\, dx
		= \frac{1}{b-a}\left[\frac{x^2}{2}\right]_{a}^{b}.
	\]
	Evaluating:
	\[
		\mathbb{E}[X] = \frac{b^2 - a^2}{2(b-a)} = \frac{a+b}{2}.
	\]
	Thus the mean of a uniform distribution is simply the midpoint.
\end{derivation}

\begin{idea}
	Symmetry around zero forces the mean to be zero without any calculation.
\end{idea}

\begin{examplebox}

	Using $a=-1$ and $b=1$:
	\[
		\mathbb{E}[X] = \frac{-1 + 1}{2} = 0.
	\]

\end{examplebox}

% -------------------------------------------------------------------

\subsection{Mean of a Random Process}


\begin{idea}
	The mean of a random process is the expected value applied
	pointwise in time. At each $t$, we average over all randomness
	that influences $X(t)$.
\end{idea}

For any process $X(t)$ constructed from random variables
and deterministic functions of $t$, linearity of expectation gives
\begin{equation}\label{eq:meanofprocess}
	m_X(t) = \mathbb{E}[X(t)].
\end{equation}

\begin{examplebox}[label=exa:uniform]
	Consider the process $X(t) = A + Bt$, where $A$ and $B$
	are independent random variables values uniform on \([-1,1]\).
	Using linearity:
	\[
		m_X(t) = \mathbb{E}[A] + t\,\mathbb{E}[B].
	\]
	For $A, B$ uniform on $[-1,1]$:
	\[
		\mathbb{E}[A] = \mathbb{E}[B] = 0,
	\]
	hence
	\[
		m_X(t) = 0.
	\]
\end{examplebox}

% -------------------------------------------------------------------

\subsection{Autocorrelation Function}
The autocorrelation function of the random process \(X(t)\) is denoted by
\(R_{XX}(t_1,t_2)\), which is often shortened to \(R_X(t_1,t_2)\) is defined as:
\begin{equation}\label{eq:autocorrelation}
	R_X(t_1,t_2) = \mathbb{E}[X(t_1)X(t_2)]
\end{equation}

From this definition we can see that the autocorrelation function can
be evaluated as:

\[
	R_X(t_1,t_2) = \int_{-\infty}^{\infty}
	\int_{-\infty}^{\infty}
	x_1 x_2 f_{X(t_1),X(t_2)}(x_1,x_2) dx_1 dx_2
\]

\begin{examplebox}
	Continuing on from \zcref{exa:uniform}. We can find the autocorrelation
	function for \(X(t) = A + Bt\) where \(A\) and \(B\) are independent random
	values uniform on \([-1,1]\).
	\begin{align*}
		R_X(t) & = \mathbb{E}[(A + Bt_1)(A + Bt_2)]                                               \\
		       & =\mathbb{E}[A^2] + \mathbb{E}[AB]t_2 + \mathbb{E}[AB]t_1 + \mathbb{E}[B^2]t_1t_2
	\end{align*}

	Because \(A\) and \(B\) are independent random variables we know that

	\[
		\mathbb{E}[AB] = \mathbb{E}[A]\mathbb{E}[B] = 0
	\]

	We can also compute \(\mathbb{E}[A^2]\) and \(\mathbb{E}[B^2]\) as follows:

	\[
		\mathbb{E}[A^2] = \int_{-1}^{1} x^2 \frac{1}{2} dx
		= \frac{1}{6} x^3 \Big|_{-1}^{1} = \frac{1}{3}
	\]

	Where we use the pdf of a uniform random variable over \([-1,1]\) which is \(\frac{1}{2}\)
	(See~\zcref{eq:uniformpdf}). Tying this all together we have:

	\[
		R_X(t_1,t_2) = \frac{1}{3} + \frac{1}{3} t_1 t_2
	\]

\end{examplebox}
